\documentclass[11pt, letterpaper]{article}

% --- Idioma y codificación ---
\usepackage[spanish]{babel}
\usepackage[utf8]{inputenc}
\usepackage[T1]{fontenc}
\usepackage{lmodern}

% --- Matemáticas ---
\usepackage{amsmath, amssymb, amsthm}

% --- Formato ---
\usepackage[margin=2.5cm]{geometry}
\usepackage{parskip}
\usepackage{booktabs}
\usepackage{hyperref}
\hypersetup{colorlinks=true, linkcolor=blue!70!black, urlcolor=blue!70!black}

% --- Entornos de teoremas ---
\theoremstyle{definition}
\newtheorem{definition}{Definición}[section]
\theoremstyle{plain}
\newtheorem{theorem}{Teorema}[section]
\newtheorem{lemma}[theorem]{Lema}
\newtheorem{corollary}[theorem]{Corolario}
\theoremstyle{remark}
\newtheorem{remark}{Observación}[section]

\title{%
  T01 --- Demostración formal:\\[4pt]
  Merge sort es $\Theta(n \log n)$
}
\author{Curso Linux--C--Rust --- B15 Estructuras de datos en C}
\date{}

\begin{document}
\maketitle
\tableofcontents

% ===================================================================
\section{La recurrencia}
% ===================================================================

El costo de merge sort se describe por:
\[
  T(n) = \begin{cases}
    \Theta(1) & \text{si } n \leq 1 \\
    2\,T(n/2) + \Theta(n) & \text{si } n > 1
  \end{cases}
\]

\begin{itemize}
  \item $2T(n/2)$: ordenar las dos mitades recursivamente.
  \item $\Theta(n)$: fusionar (merge) las dos mitades ordenadas.
\end{itemize}

Para simplificar, asumimos $n = 2^k$ (potencia de~2). El resultado se
extiende a todo $n$ con pisos y techos, sin cambiar la clase asintótica.

Usamos $T(n) = 2T(n/2) + cn$ con $c > 0$ constante, y $T(1) = d$ con
$d > 0$ constante.

% ===================================================================
\section{Demostración 1: árbol de recursión}
% ===================================================================

\begin{theorem}\label{thm:rec-tree}
$T(n) = \Theta(n \log n)$.
\end{theorem}

\begin{proof}
Expandimos la recurrencia en forma de árbol. Cada nodo representa una
llamada, con su costo de merge (sin contar las llamadas recursivas):

\begin{center}
\begin{tabular}{lcl}
Nivel 0: & $cn$ & $\to cn$ \\
Nivel 1: & $2 \times cn/2$ & $\to cn$ \\
Nivel 2: & $4 \times cn/4$ & $\to cn$ \\
$\vdots$ & $\vdots$ & $\vdots$ \\
Nivel $j$: & $2^j \times cn/2^j$ & $\to cn$ \\
$\vdots$ & $\vdots$ & $\vdots$ \\
Nivel $k$: & $n \times c$ & $\to cn$ \\
\end{tabular}
\end{center}

\textbf{Número de niveles.} La recursión divide $n$ por~2 en cada
nivel. Llega al caso base ($n = 1$) en el nivel $k = \log_2 n$. Hay
$\log_2 n + 1$ niveles en total (nivel~0 a nivel~$\log_2 n$).

\textbf{Trabajo por nivel.} En el nivel~$j$ hay $2^j$ subproblemas,
cada uno de tamaño $n/2^j$, con costo de merge $c \cdot n/2^j$. El
trabajo total en el nivel~$j$ es:
\[
  2^j \cdot \frac{cn}{2^j} = cn
\]
Cada nivel contribuye exactamente $cn$ de trabajo, independientemente
del nivel.

\textbf{Trabajo total.} Sumando sobre todos los niveles:
\[
  T(n) = \sum_{j=0}^{\log_2 n} cn
       = cn \cdot (\log_2 n + 1)
       = cn\log_2 n + cn
       = \Theta(n \log n) \qedhere
\]
\end{proof}

% ===================================================================
\section{Demostración 2: Teorema Maestro}
% ===================================================================

\begin{theorem}[Teorema Maestro]\label{thm:master}
Sea $T(n) = aT(n/b) + f(n)$ con $a \geq 1$, $b > 1$, y $f(n)$
asintóticamente positiva. Sea $c_{\mathrm{crit}} = \log_b a$. Entonces:

\begin{itemize}
  \item \textbf{Caso 1.} Si
    $f(n) = O(n^{c_{\mathrm{crit}} - \varepsilon})$ para algún
    $\varepsilon > 0$, entonces
    $T(n) = \Theta(n^{c_{\mathrm{crit}}})$.
  \item \textbf{Caso 2.} Si
    $f(n) = \Theta(n^{c_{\mathrm{crit}}} \log^k n)$ para algún
    $k \geq 0$, entonces
    $T(n) = \Theta(n^{c_{\mathrm{crit}}} \log^{k+1} n)$.
  \item \textbf{Caso 3.} Si
    $f(n) = \Omega(n^{c_{\mathrm{crit}} + \varepsilon})$ para algún
    $\varepsilon > 0$ y $af(n/b) \leq cf(n)$ para algún $c < 1$,
    entonces $T(n) = \Theta(f(n))$.
\end{itemize}
\end{theorem}

\begin{theorem}\label{thm:master-app}
La recurrencia de merge sort $T(n) = 2T(n/2) + cn$ satisface
$T(n) = \Theta(n \log n)$.
\end{theorem}

\begin{proof}
Identificamos parámetros:
\begin{itemize}
  \item $a = 2$ (dos subproblemas).
  \item $b = 2$ (cada uno de tamaño $n/2$).
  \item $f(n) = cn = \Theta(n)$.
\end{itemize}

Calculamos $c_{\mathrm{crit}}$:
\[
  c_{\mathrm{crit}} = \log_b a = \log_2 2 = 1
\]

Comparamos $f(n)$ con $n^{c_{\mathrm{crit}}} = n$:
\[
  f(n) = cn = \Theta(n) = \Theta(n^{c_{\mathrm{crit}}} \cdot \log^0 n)
\]

Esto es el \textbf{Caso~2} con $k = 0$. Por el Teorema Maestro:
\[
  T(n) = \Theta(n^{c_{\mathrm{crit}}} \log^{k+1} n)
       = \Theta(n^1 \cdot \log^1 n)
       = \Theta(n \log n) \qedhere
\]
\end{proof}

% ===================================================================
\section{Demostración 3: verificación por sustitución (inducción)}
% ===================================================================

\begin{theorem}\label{thm:substitution}
Si $T(1) = d$ y $T(n) = 2T(n/2) + cn$ para $n \geq 2$ (con $n$
potencia de~2), entonces $T(n) = cn\log_2 n + dn$.
\end{theorem}

\begin{proof}
Por inducción sobre $k$ donde $n = 2^k$.

\textbf{Base ($k = 0$, $n = 1$):}
\[
  T(1) = d = c \cdot 1 \cdot \log_2 1 + d \cdot 1 = 0 + d = d
  \quad \checkmark
\]

\textbf{Paso inductivo.} Sea $n = 2^k$ con $k \geq 1$. Supongamos que
$T(n/2) = c(n/2)\log_2(n/2) + d(n/2)$ (hipótesis inductiva). Entonces:
\begin{align*}
  T(n) &= 2T(n/2) + cn \\
       &= 2\left[\frac{cn}{2}\log_2\frac{n}{2}
         + \frac{dn}{2}\right] + cn \\
       &= cn\log_2\frac{n}{2} + dn + cn \\
       &= cn(\log_2 n - 1) + dn + cn \\
       &= cn\log_2 n - cn + dn + cn \\
       &= cn\log_2 n + dn \qedhere
\end{align*}
\end{proof}

% ===================================================================
\section{Extensión a $n$ general (no potencia de 2)}
% ===================================================================

Cuando $n$ no es potencia de~2, la recurrencia exacta es:
\[
  T(n) = T\!\left(\left\lfloor \frac{n}{2} \right\rfloor\right)
       + T\!\left(\left\lceil \frac{n}{2} \right\rceil\right) + cn
\]

\begin{theorem}\label{thm:general-n}
La clase asintótica no cambia: $T(n) = \Theta(n \log n)$.
\end{theorem}

\begin{proof}
Sea $n' = 2^{\lceil \log_2 n \rceil}$ la potencia de~2 inmediatamente
$\geq n$. Entonces $n \leq n' < 2n$, y:
\[
  T(n) \leq T(n') = cn'\log_2 n' + dn'
  \leq c(2n)\log_2(2n) + d(2n) = O(n\log n)
\]

La cota inferior se obtiene análogamente con
$n'' = 2^{\lfloor \log_2 n \rfloor} \geq n/2$. Combinando:
$T(n) = \Theta(n \log n)$.
\end{proof}

% ===================================================================
\section{Demostración 4: merge sort es
  \texorpdfstring{$\Theta(n \log n)$}{Theta(n log n)} en todos los casos}
% ===================================================================

\begin{theorem}\label{thm:all-cases}
Merge sort tiene complejidad $\Theta(n \log n)$ en el mejor, promedio
y peor caso.
\end{theorem}

\begin{proof}
La estructura recursiva es \textbf{independiente de la entrada}: merge
sort siempre divide a la mitad y siempre fusiona. Lo único que varía es
el número de comparaciones dentro de cada merge.

\emph{Peor caso.} El merge de dos subarrays de tamaños $p$ y $q$
ejecuta a lo sumo $p + q - 1$ comparaciones (cuando los elementos se
intercalan perfectamente). En cada nivel, las comparaciones totales son
a lo sumo $n - (\text{número de subarrays}) \leq n$. Con $\log_2 n$
niveles: $T_{\text{worst}} \leq n \log_2 n = O(n \log n)$.

\emph{Mejor caso.} Cada merge ejecuta al menos $\min(p, q)$
comparaciones (una mitad se agota antes de tocar la otra). Para mitades
iguales ($p = q = n/2^j$), esto es $n/2^j$. En cada nivel, las
comparaciones totales son al menos $n/2$. Con $\log_2 n$ niveles:
$T_{\text{best}} \geq (n/2)\log_2 n = \Omega(n \log n)$.

Combinando: merge sort es $\Theta(n \log n)$ en todos los casos.
\end{proof}

% ===================================================================
\section{Resumen de resultados}
% ===================================================================

\begin{table}[h]
\centering
\begin{tabular}{llp{5.5cm}}
\toprule
\textbf{Resultado} & \textbf{Método} & \textbf{Clave} \\
\midrule
$T(n) = cn\log_2 n + dn$
  & Árbol de recursión
  & $cn$ trabajo por nivel $\times$ $\log_2 n + 1$ niveles \\
$T(n) = \Theta(n \log n)$
  & Teorema Maestro (caso~2)
  & $a = b = 2$, $f(n) = \Theta(n^{\log_2 2})$ \\
Fórmula exacta verificada
  & Inducción (sustitución)
  & $cn(\log_2 n {-} 1) + dn + cn = cn\log_2 n + dn$ \\
$n$ general no cambia la clase
  & Redondeo a potencia de~2
  & $n \leq 2^{\lceil \log_2 n \rceil} < 2n$ \\
$\Theta(n \log n)$ en todos los casos
  & Cotas de merge
  & Mejor: $\geq n/2$ por nivel; peor: $\leq n$ por nivel \\
\bottomrule
\end{tabular}
\end{table}

\end{document}
